% R599: whole EGA I Chapter0 3.5.5.
% OpenAI Codex, GPT-6.1 Sol, Ultra; no human review claimed.
\begin{env}[3.5.5]
\label{0.3.5.5-fr}
Sea $\psi':Y\to Z$ una aplicación continua, y supongamos que todo
prehaz $\mathcal{H}$ sobre $Z$ con valores en $\boldsymbol{K}$ admita una
imagen inversa $\psi'^*(\mathcal{H})$ por $\psi'$. Entonces (con las
hipótesis de (3.5.4)) todo prehaz $\mathcal{H}$ sobre $Z$ con valores en
$\boldsymbol{K}$ admite una imagen inversa por $\psi''=\psi'\circ\psi$ y se tiene
un isomorfismo canónico funtorial
\begin{equation*}
  \psi''^*(\mathcal{H})\xrightarrow{\sim}
  \psi^*(\psi'^*(\mathcal{H}))
\tag{3.5.5.1}\label{0.3.5.5.1-fr}
\end{equation*}

\oldpage[0\textsubscript{I}]{33}

Esto se deduce, en efecto, inmediatamente de las definiciones, teniendo en cuenta que
$\psi''_*=\psi'_*\circ\psi_*$. Además, si
$u:\mathcal{G}\to\psi_*(\mathcal{F})$ es un $\psi$-morfismo,
$v:\mathcal{H}\to\psi'_*(\mathcal{G})$ un $\psi'$-morfismo, y
$w=\psi'_*(u)\circ v$ su compuesto (3.5.2), se comprueba inmediatamente que
$w^\sharp$ es el morfismo compuesto
\[
  w^\sharp:\psi^*(\psi'^*(\mathcal{H}))
  \xrightarrow{\psi^*(v^\sharp)}
  \psi^*(\mathcal{G})
  \xrightarrow{u^\sharp}\mathcal{F}.
\]
\end{env}
